\appendix

\chapter{Captures d'écran de la version déployée}
\label{ch:captures}

Captures récentes de la plateforme en production \url{https://lexhub.space}
(espace cabinet Démo Atlas). Elles illustrent le parcours public, l'espace cabinet,
les documents, l'atelier plaidoirie, la collaboration d'équipe et l'administration
plateforme.

\section*{Espace public}

\begin{figure}[H]
\centering
\includegraphics[width=0.92\textwidth]{01-accueil.png}
\caption{Accueil LexHub (page publique)}
\label{fig:accueil}
\end{figure}

\begin{figure}[H]
\centering
\includegraphics[width=0.92\textwidth]{02-connexion.png}
\caption{Connexion sécurisée}
\label{fig:connexion}
\end{figure}

\section*{Espace cabinet --- vue d'ensemble}

\begin{figure}[H]
\centering
\includegraphics[width=0.92\textwidth]{10-tableau-de-bord.png}
\caption{Tableau de bord --- dossiers, documents, essai et crédits}
\label{fig:dashboard}
\end{figure}

\begin{figure}[H]
\centering
\includegraphics[width=0.92\textwidth]{11-dossiers.png}
\caption{Liste des dossiers du cabinet}
\label{fig:dossiers}
\end{figure}

\begin{figure}[H]
\centering
\includegraphics[width=0.92\textwidth]{24-dossier-detail.png}
\caption{Fiche dossier --- identité des parties et plaidoirie liée}
\label{fig:casedetail}
\end{figure}

\section*{Documents}

\begin{figure}[H]
\centering
\includegraphics[width=0.92\textwidth]{13-documents.png}
\caption{Espace Documents --- dépôt et statut d'extraction}
\label{fig:documents}
\end{figure}

\begin{figure}[H]
\centering
\includegraphics[width=0.92\textwidth]{13b-document-detail.png}
\caption{Détail document --- extrait, synthèse IA et questions}
\label{fig:docdetail}
\end{figure}

\section*{Cœur métier --- plaidoirie}

\begin{figure}[H]
\centering
\includegraphics[width=0.92\textwidth]{14b-plaidoirie-workspace.png}
\caption{Atelier plaidoirie --- étape 1 (dossier \& pièces)}
\label{fig:plaidoirie}
\end{figure}

\begin{figure}[H]
\centering
\includegraphics[width=0.92\textwidth]{14c-transcription-fusha.png}
\caption{Transcription / correction (arabe fusḥā)}
\label{fig:transcription}
\end{figure}

\begin{figure}[H]
\centering
\includegraphics[width=0.92\textwidth]{19-export-document.png}
\caption{Export de la plaidoirie (Word .docx / PDF)}
\label{fig:export}
\end{figure}

\section*{Équipe et activité}

\begin{figure}[H]
\centering
\includegraphics[width=0.92\textwidth]{16-equipe.png}
\caption{Équipe --- invitations et rôles}
\label{fig:equipe}
\end{figure}

\begin{figure}[H]
\centering
\includegraphics[width=0.92\textwidth]{23-activite-cabinet.png}
\caption{Journal d'activité du cabinet}
\label{fig:activite}
\end{figure}

\section*{Administration plateforme}

\begin{figure}[H]
\centering
\includegraphics[width=0.92\textwidth]{22a-admin-dashboard.png}
\caption{Admin plateforme --- pilotage LexHub}
\label{fig:admin}
\end{figure}

\section{Sources des diagrammes UML}
\label{sec:sources-uml}

Les diagrammes UML du rapport sont ancrés sur le code de production. Les sources éditables PlantUML (et Mermaid pour l'ER / UC) se trouvent dans \texttt{docs/rapport-pfa/diagrams/} :

\begin{itemize}[leftmargin=1.4em]
  \item \texttt{seq-register-otp.puml} (\texttt{seq-otp.puml}) --- inscription OTP ;
  \item \texttt{seq-login-turnstile.puml} --- login + Turnstile ;
  \item \texttt{seq-dictation.puml} --- dictée $\rightarrow$ clarification / composition / révision ;
  \item \texttt{seq-validate-assign.puml} --- validation, assignation, export Word/PDF ;
  \item \texttt{seq-document-ocr.puml} --- dépôt document et OCR ;
  \item \texttt{seq-billing.puml} --- plans, packs, factures payées ;
  \item \texttt{activity-plea.puml} --- activité atelier ;
  \item \texttt{state-plaidoirie.puml} --- états \texttt{draft} / \texttt{validated} / \texttt{finalized} ;
  \item \texttt{class-lexhub.puml}, \texttt{er-lexhub.mmd} --- extrait de classes / ER ;
  \item \texttt{deploy-lexhub.puml} --- déploiement ;
  \item \texttt{uc-lexhub.puml}, \texttt{uc-lexhub.mmd} --- cas d'utilisation.
\end{itemize}

Le chapitre conception reprend ces flux en TikZ ; les fichiers \texttt{.puml} restent la référence éditable.
